
\documentclass[11pt,a4paper]{article}
\usepackage[a4paper,margin=2.1cm]{geometry}
\usepackage[spanish,es-noquoting]{babel}
\usepackage[T1]{fontenc}
\usepackage[utf8]{inputenc}
\usepackage{lmodern,microtype}
\usepackage{amsmath,amssymb,amsthm,mathtools}
\usepackage{booktabs,longtable,array,enumitem}
\usepackage{graphicx}
\usepackage{hyperref}
\hypersetup{colorlinks=true,linkcolor=blue,urlcolor=blue,citecolor=blue}

\newtheorem{teorema}{Teorema}
\newtheorem{proposicion}[teorema]{Proposición}
\newtheorem{definicion}[teorema]{Definición}
\newtheorem{observacion}[teorema]{Observación}

\newcommand{\W}{W_{12}}
\newcommand{\muW}{\mu_W}
\newcommand{\sigmaB}{\sigma_B}
\newcommand{\BC}{B_{\mathfrak C}}
\newcommand{\Halpha}{H_\alpha^-}
\newcommand{\Hplus}{H_\alpha^+}
\newcommand{\He}{H_e}
\newcommand{\Hphi}{H_\varphi}
\newcommand{\Happ}{H_{\rm APP}}
\newcommand{\Ppi}{P_\pi}

\title{\bfseries Comunicado HMT: holonomía Mathieu del campo de estrellas de Witt\\
\large Involuciones locales, generación de \(M_{12}\), rigidez del marco HMT y lectura de \(\alpha\)}
\author{Ventana inter-chats}
\date{}

\begin{document}
\maketitle

\begin{abstract}
Se verifica que el campo de emparejamientos de Witt no sólo organiza las hexadas incidentes a cada 4-cara, sino que genera la simetría global del sistema. A cada 4-cara \(B\subset[12]\) se asocia una involución \(\sigma_B\), que fija \(B\) y transpone los cuatro pares de \(\mu_W(B)\). Por enumeración exacta sobre las 132 hexadas de \(W_{12}=S(5,6,12)\), cada \(\sigma_B\) preserva \(W_{12}\), y el grupo generado por todas ellas tiene orden \(95040\), por tanto coincide con \(M_{12}\). Además, la configuración HMT
\[
(\Halpha,\He,\Hphi,\Happ)
\]
tiene estabilizador trivial. Esto muestra que el borde HMT no sólo vive en \(W_{12}\): rigidifica por completo su gauge de Mathieu.
\end{abstract}

\section{Campo de emparejamientos y sus involuciones}

Para cada 4-cara \(B\subset[12]\), el diseño de Witt determina cuatro hexadas incidentes. Los dos puntos añadidos a \(B\) en cada hexada forman cuatro pares disjuntos:
\[
\mu_W(B)=\{H\setminus B:\ H\in W_{12},\ B\subset H\}.
\]
Estos cuatro pares particionan \([12]\setminus B\). Definimos la involución local
\[
\sigma_B
\]
como la permutación que fija todos los puntos de \(B\) y transpone cada uno de los cuatro pares de \(\mu_W(B)\).

\begin{proposicion}[certificado finito]
Para toda 4-cara \(B\), la involución \(\sigma_B\) preserva el diseño de Witt:
\[
\sigma_B(W_{12})=W_{12}.
\]
\end{proposicion}

\begin{observacion}
Esta proposición se verifica por enumeración exacta de las 132 hexadas. No depende de \(\alpha\), ni de metrología, ni de la lectura archimedeana.
\end{observacion}

\section{Generación de \(M_{12}\)}

Sea
\[
\Gamma_W=\langle \sigma_B:\ B\in\binom{[12]}4\rangle.
\]
La enumeración de Schreier--Sims da
\[
|\Gamma_W|=95040.
\]
Como \(95040\) es el orden de \(\operatorname{Aut}(W_{12})\cong M_{12}\), se obtiene:

\begin{teorema}[holonomía Mathieu del campo de estrellas]
Las involuciones locales del campo de emparejamientos generan la simetría completa de Witt:
\[
\boxed{
\langle \sigma_B:\ |B|=4\rangle=M_{12}.
}
\]
\end{teorema}

\begin{observacion}
Esta es la forma finita de una holonomía: los datos locales de emparejamiento, al componerse, generan la simetría global del borde.
\end{observacion}

\section{La estrella de \(\alpha\)}

La 4-cara alta del corrector es
\[
\BC=\{4,6,9,10\}.
\]
Su campo de emparejamientos es
\[
\mu_W(\BC)=
\{\{1,3\},\{2,11\},\{8,12\},\{5,7\}\}.
\]
Por tanto:
\[
\sigma_{\BC}=(1\,3)(2\,11)(5\,7)(8\,12),
\]
fijando
\[
4,6,9,10.
\]
La partición de signo del cierre de \(\alpha\) es
\[
\Halpha=\{4,5,6,7,9,10\},
\]
\[
\Hplus=\{1,2,3,8,11,12\}.
\]
Así, los pares
\[
\{1,3\},\quad \{2,11\},\quad \{8,12\}
\]
son positivos, y
\[
\{5,7\}
\]
es el único par negativo. En consecuencia:
\[
\Halpha=\BC\cup\{5,7\}.
\]

\section{Rigidez del marco HMT}

La cadena de estabilizadores dentro de \(M_{12}\) es:

\[
\begin{array}{c|c}
\text{configuración marcada} & \text{orden del estabilizador}\\
\hline
\Halpha & 720\\
(\Halpha,6) & 120\\
\BC & 192\\
(\BC,\Halpha) & 48\\
(\BC,\Halpha,\He) & 16\\
(\Halpha,\He,\Hphi) & 2\\
(\Halpha,\He,\Hphi,\Happ) & 1
\end{array}
\]

\begin{teorema}[rigidez del marco HMT]
La configuración
\[
(\Halpha,\He,\Hphi,\Happ)
\]
tiene estabilizador trivial en \(M_{12}\):
\[
\operatorname{Stab}_{M_{12}}(\Halpha,\He,\Hphi,\Happ)=\{1\}.
\]
\end{teorema}

\begin{observacion}
Esto implica que las incidencias HMT no sólo seleccionan objetos dentro de \(W_{12}\). Seleccionan un marco completo y eliminan todo gauge residual de \(M_{12}\).
\end{observacion}

\section{No circularidad}

El núcleo de la estrella de \(\alpha\) se obtiene por dos rutas independientes:
\[
\BC=\{i:K_i\ge729\}
\]
desde la anatomía \(729+271\) del corrector, y
\[
\BC=\He\cap\Ppi
\]
desde la lectura Golay de los canales \(e\) y \(\pi\).

El par negativo \(\{5,7\}\) se obtiene del signo del cierre, y el lift real de la cochain normalizada produce \(\alpha\). La cadena no usa \(\alpha\) para construir \(B_{\mathfrak C}\), ni usa \(M_{12}\) para ajustar el valor real.

\section{Interpretación}

La lectura refinada es:
\[
\boxed{
\alpha \text{ es el lift archimedeano del pétalo negativo de una estrella de Witt cuya 4-cara nuclear es la corona }729\text{ del corrector.}
}
\]

Pero ahora se añade una capa más:
\[
\boxed{
\text{las estrellas locales de Witt generan }M_{12},
\text{ y la configuración HMT fija completamente ese marco.}
}
\]

Por tanto, HMT no sólo toca \(W_{12}\); selecciona una estrella polarizada dentro de un campo cuya holonomía finita es Mathieu.

\section{Conclusión}

La morfología final de esta etapa es:
\[
\text{microtransporte}
\longrightarrow
\text{4-cara de interfase}
\longrightarrow
\text{estrella de Witt}
\longrightarrow
\text{pétalo negativo}
\longrightarrow
\text{lift real } \alpha
\]
y, en paralelo,
\[
\text{campo de estrellas}
\longrightarrow
M_{12}
\longrightarrow
\text{Golay--Leech--Moonshine}.
\]

La frase de cierre es:
\[
\boxed{
\text{el número deja de ser valor y se vuelve sección de una geometría finita de borde.}
}
\]

\end{document}
